% 第 4 章 Tolerance for Incomplete or Mismatched Data
\bichapter{对不完整或不匹配数据的容错}{Tolerance for Incomplete or Mismatched Data}
\label{chap:tolerance}

即使 RTL netlist 与 library 之间、logic library 与 physical library 之间，或 top-level RTL instantiation 与 RTL subblock 的 interface pin 定义之间存在不完整或不匹配的数据，DC Explorer 仍可继续执行 link process。
对于这些数据不一致，工具会汇总 link error 和 warning message，并提供详细报告。
当工具遇到 infeasible path 时，会推断 path assumption，并继续优化其他 critical path。

对不完整或不匹配数据的容错默认启用。在 DC Explorer 读入并链接设计之前，可以更改控制特定容错行为的部分设置，例如 bus 与 bit-blasted 命名风格。

要进一步了解对不完整或不匹配数据的容错，请参阅：
\begin{itemize}
  \item 容错类别
  \item 默认容错设置
  \item 调整总线引脚命名风格的容错设置
  \item 解析总线与 bit-blasted 命名风格
  \item 允许引脚名同义词
  \item 允许未连接的接口引脚
  \item 在 netlist 中包含 anchor cell
  \item 创建报告与缺失约束
\end{itemize}

% ============================================================
\section{容错类别}
\subsection*{Tolerance Categories}

表~\ref{tab:tol-cats} 列出 DC Explorer 支持的对不完整或不匹配数据的容错类别。
勾选（X）表示该容错类别适用且受支持。

\begin{table}[htbp]
\centering
\caption{对不完整或不匹配数据的容错类别}
\label{tab:tol-cats}
\small
\begin{tabularx}{\textwidth}{@{}X ccc@{}}
\toprule
\textbf{容错类别} &
\textbf{Library 定义 vs instantiation} &
\textbf{Top-level RTL vs RTL block} &
\textbf{Logic library vs physical library} \\
\midrule
不匹配的总线引脚命名风格 & X & X & X \\
总线 vs bit-blasted & X & X & X \\
引脚名大小写不匹配 & X & X & X \\
逻辑库中缺失单元 & X & & \\
物理库中缺失单元 & & & X \\
逻辑库中缺失引脚 & X & & X \\
物理库中缺失引脚 & & & X \\
RTL 模块中缺失引脚 & & X & \\
缺失 RTL 模块 & & X & \\
引脚方向不匹配 & & & X \\
引脚名同义词（引脚映射） & X & X & X \\
端口宽度不匹配 & X & X & \\
未连接的接口引脚 & X & X & \\
\bottomrule
\end{tabularx}
\end{table}

\begin{seeAlsoBox}
\begin{itemize}
  \item 默认容错设置
  \item 调整总线引脚命名风格的容错设置
  \item 解析总线与 bit-blasted 命名风格
  \item 允许引脚名同义词
  \item 允许未连接的接口引脚
  \item 报告不完整或不匹配数据
\end{itemize}
\end{seeAlsoBox}

\subsection{限制}
\subsubsection*{Limitations}

DC Explorer 支持 Verilog 文件中所有容错类别的不匹配，但不支持 VHDL 文件中 subblock instantiation 与 top-level component instantiation 之间的不匹配。
以下 VHDL 代码属于「RTL module 中缺失 pin」这一容错类别。
代码显示 SUB1 与 TOP 两个 instantiation 之间存在不匹配，其中 SUB1 instantiation 缺少 MSGPIN pin。
分析设计时，工具会对此类不匹配发出 error message。
\begin{lstlisting}
entity SUB1 is
   port (AA, BB: in bit; CC: out bit);
end entity SUB1;

architecture RTL of SUB1 is
begin
   CC <= AA and BB;
end architecture RTL;

library WORK;
use WORK.all;

entity TOP is
   port (a, b, e: in bit; c: out bit);
end entity TOP;

architecture RTL of TOP is
begin
   U1 : entity WORK.SUB1(RTL)
   port map (AA=>a, BB =>b, MSGPIN=>e, CC=>c);
end architecture RTL;
\end{lstlisting}

\begin{seeAlsoBox}
容错类别（上文）。
\end{seeAlsoBox}

% ============================================================
\section{默认容错设置}
\subsection*{Default Tolerance Setup}

默认情况下，DC Explorer 已针对下列容错类别完成设置；除 bus pin naming style 的容错外，其他设置均不可配置。
有关 DC Explorer 的默认行为，请参阅各容错类别的说明：
\begin{itemize}
  \item 总线引脚命名风格
  \item 引脚名大小写不匹配
  \item 逻辑库中缺失单元
  \item 物理库中缺失单元
  \item 逻辑库中缺失引脚
  \item FRAM 视图中缺失引脚
  \item 逻辑库与物理库之间的引脚方向不匹配
  \item 端口宽度不匹配
\end{itemize}

\subsection{总线引脚命名风格}
\subsubsection*{Bus Pin Naming Styles}

默认情况下，DC Explorer 可根据 \varname{bit\_blasted\_bus\_linking\_naming\_styles} 变量定义的命名约定，将 RTL block 与 logic library 中的对应 block 进行匹配。
要为不同总线命名风格调整该变量的默认设置，见「调整总线引脚命名风格的容错设置」。
表~\ref{tab:bus-match-default} 给出基于该变量默认设置的匹配示例。

\begin{table}[htbp]
\centering
\caption{默认总线名匹配示例}
\label{tab:bus-match-default}
\begin{tabularx}{\textwidth}{@{}X X@{}}
\toprule
\textbf{RTL} & \textbf{逻辑库文件} \\
\midrule
\texttt{OPRNDA\_0\_} & \texttt{OPRNDA[0]} \\
\texttt{OPRNDA\_1\_} & \texttt{OPRNDA[1]} \\
\texttt{OPRNDB(0)} & \texttt{OPRNDB\_0\_} \\
\texttt{OPRNDB(1)} & \texttt{OPRNDB\_1\_} \\
\texttt{OPRNDC[0]} & \texttt{OPRNDC(0)} \\
\texttt{OPRNDC[1]} & \texttt{OPRNDC(1)} \\
\bottomrule
\end{tabularx}
\end{table}

\begin{seeAlsoBox}
调整总线引脚命名风格的容错设置。
\end{seeAlsoBox}

\subsection{引脚名大小写不匹配}
\subsubsection*{Case Mismatches Between Pin Names}

默认情况下，DC Explorer 允许下列对象之间存在 pin name 大小写不匹配：
\begin{itemize}
  \item 顶层例化与 RTL 块
  \item 逻辑库与 RTL 例化
  \item 逻辑库与物理库（FRAM 视图）
\end{itemize}

图~\ref{fig:case-mismatch} 给出顶层例化与 RTL 块之间引脚名大小写不匹配的示例。加载到内存中的设计显示这些不匹配已被解析。

\figplaceholder{Figure 4: Example of Case Mismatches Between Pin Name}{引脚名大小写不匹配示例}{fig:case-mismatch}

\subsection{逻辑库中缺失单元}
\subsubsection*{Missing Cells in Logic Libraries}

在 link process 中，即使 logic library 中缺少 cell reference，DC Explorer 仍可完成 cell linking，如图~\ref{fig:miss-cell-logic} 所示。
在此类情况下，DC Explorer 不发出 error message，但解析后得到的 library cell、macro 或 module 不包含 timing arc。

\figplaceholder{Figure 5: Tolerance for Missing Cells in Logic Libraries}{对逻辑库中缺失单元的容错}{fig:miss-cell-logic}

\subsection{物理库中缺失单元}
\subsubsection*{Missing Cells in Physical Libraries}

当 physical library 中缺少某个 cell 时，DC Explorer 执行下列操作以解决该不匹配：
\begin{itemize}
  \item 在 local Milkyway design library 中为缺失 cell 创建 FRAM view，并使用 one-unit-tile bounding box
  \item 从 logic library 推导缺失 cell 的 pin direction
  \item 在 FRAM view 的原点添加 physical pin
\end{itemize}

对物理库中较大的缺失宏，应按黑盒流程中所述步骤解析不匹配。
有关黑盒流程的更多信息，见 \textit{IC Compiler Design Planning User Guide} 与 \textit{IC Compiler User Guide}。
图~\ref{fig:miss-cell-phys} 给出为子逻辑单元创建的 FRAM 视图。

\figplaceholder{Figure 6: Examples of Missing Cells in Physical Libraries}{物理库中缺失单元的示例}{fig:miss-cell-phys}

\subsection{逻辑库中缺失引脚}
\subsubsection*{Missing Pins in Logic Libraries}

当 logic library cell 或 RTL reference 存在缺失 pin 时，DC Explorer 通过下列方式容许该不匹配：
\begin{itemize}
  \item 为每个缺失 pin 创建 anchor cell，以保留连接并防止该 pin 被移除
  \item 忽略缺失 pin 上的 timing constraint
  \item 写出设计时，自动从 ASCII netlist 中移除 anchor cell，并恢复缺失 pin 的连接
\end{itemize}

如图~\ref{fig:miss-pin-logic} 所示，为缺失的 \texttt{wr} 与 \texttt{ovf} 引脚创建了两个锚点单元。

\figplaceholder{Figure 7: Example of Missing Pins in Logic Libraries}{逻辑库中缺失引脚的示例}{fig:miss-pin-logic}

\begin{seeAlsoBox}
在网表中包含锚点单元。
\end{seeAlsoBox}

\subsection{FRAM 视图中缺失引脚}
\subsubsection*{Missing Pins in FRAM Views}

当 logic library cell 的 pin 数量多于 FRAM view 时，DC Explorer 会为每个额外 pin 创建一个新的 physical pin。
新的 physical pin 添加在 cell bounding box 的左下角，并放置在最低的允许 routing layer 上。

如图~\ref{fig:miss-pin-fram} 所示，为 FRAM 视图中缺失的 \texttt{wr} 引脚创建了虚设物理引脚，并放置在子块的左下角。

\figplaceholder{Figure 8: Example of a Missing Pin in the FRAM View}{FRAM 视图中缺失引脚的示例}{fig:miss-pin-fram}

\subsection{逻辑库与物理库之间的引脚方向不匹配}
\subsubsection*{Pin Direction Mismatches Between Logic and Physical Libraries}

DC Explorer 允许 port direction mismatch。例如，当 logic library 中指定的 pin direction 与 FRAM view 中的方向不同时，DC Explorer 使用 logic library 中指定的方向，为 logic library cell 创建 physical link。
图~\ref{fig:pin-dir-mismatch} 说明 DC Explorer 如何自动解析 \texttt{di} 与 \texttt{ovf} 引脚方向不匹配。

\figplaceholder{Figure 9: Example of Pin Direction Mismatches}{引脚方向不匹配示例}{fig:pin-dir-mismatch}

\subsection{端口宽度不匹配}
\subsubsection*{Port Width Mismatches}

DC Explorer 允许 port width mismatch。例如，当 top-level design 所实例化的 RTL block 存在 port width mismatch 时，DC Explorer 会截断或 zero-extend bus，并添加 anchor cell，以保留 top level 上的额外 signal connection，如图~\ref{fig:port-width} 所示。

\figplaceholder{Figure 10: Example of Port Width Mismatches}{端口宽度不匹配示例}{fig:port-width}

% ============================================================
\section{调整总线引脚命名风格的容错设置}
\subsection*{Adjusting the Tolerance Setting for Bus Pin Naming Styles}

DC Explorer 默认可匹配某些总线引脚命名风格。
通过设置 \varname{bit\_blasted\_bus\_linking\_naming\_styles} 变量，可使 DC Explorer 匹配其他总线引脚命名风格。
该变量默认为：
\begin{lstlisting}
"%s\[%d\] %s(%d) %s_%d_"
\end{lstlisting}
要允许 DC Explorer 匹配表~\ref{tab:bus-styles} 所示的总线引脚命名风格，输入：
\begin{lstlisting}
de_shell> set_app_var enable_bit_blasted_bus_linking true
de_shell> set_app_var bit_blasted_bus_linking_naming_styles \
   "%s_%d %s\[%d\] %s<%d>"
\end{lstlisting}

在此示例中，\verb|%s_%d| 表示以一个下划线和一个数字结尾的字符串。
该形式的名称可匹配由同一字符串后跟方括号内同一数字（\verb|%s\[%d\]|），或同一字符串后跟尖括号内同一数字（\verb|%s<%d>|）组成的任意名称。
在匹配时，采用这三种指定形式的任意一对名称均视为等价。
表~\ref{tab:bus-styles} 给出此示例的若干匹配名称。

\begin{table}[htbp]
\centering
\caption{RTL 与逻辑库文件之间不匹配的总线引脚命名风格}
\label{tab:bus-styles}
\begin{tabularx}{\textwidth}{@{}X X@{}}
\toprule
\textbf{RTL} & \textbf{逻辑库文件} \\
\midrule
\texttt{OPRNDA\_0} & \texttt{OPRNDA[0]} \\
\texttt{OPRNDA\_1} & \texttt{OPRNDA[1]} \\
\texttt{OPRNDB<0>} & \texttt{OPRNDB\_0} \\
\texttt{OPRNDB<1>} & \texttt{OPRNDB\_1} \\
\texttt{OPRNDC<0>} & \texttt{OPRNDC[0]} \\
\texttt{OPRNDA\_0} & \texttt{OPRNDA[0]} \\
\bottomrule
\end{tabularx}
\end{table}

\begin{seeAlsoBox}
总线引脚命名风格。
\end{seeAlsoBox}

% ============================================================
\section{解析总线与 bit-blasted 命名风格}
\subsection*{Resolving Bus and Bit-Blasted Naming Styles}

要使 DC Explorer 匹配不同的总线与 bit-blasted 命名风格：
\begin{enumerate}
  \item 将 \varname{enable\_bit\_blasted\_bus\_linking} 变量设为 \texttt{true}（默认值）。
  \item 将 \varname{bit\_blasted\_bus\_linking\_naming\_styles} 变量设为具体命名风格。
\end{enumerate}

如图~\ref{fig:bus-blasted} 所示，子 RTL 块的 \verb|\sel[1]| 与 \verb|\sel[0]| bit-blasted 信号不匹配。可设置下列变量以解析不匹配：
\begin{lstlisting}
de_shell> set_app_var enable_bit_blasted_bus_linking true
de_shell> set_app_var bit_blasted_bus_linking_naming_styles \
   "%s\[%d\] %s(%d) %s_%d_"
\end{lstlisting}

\figplaceholder{Figure 11: Example of Bus and Bus-Blasted Naming Styles}{总线与 bit-blasted 命名风格示例}{fig:bus-blasted}

使用与上一示例相同的两个变量设置，DC Explorer 可匹配多维总线，例如下列 RTL 示例中 \texttt{sub1} RTL 块的 \texttt{a} 总线与 \texttt{b} 总线：
\begin{lstlisting}
module sub1 (input \a[1][1], \a[1][0], \a[0][1], \a[0][0],
             input [1:0][1:0]b,
             output [1:0][1:0] c);

module top (input [1:0][1:0] ain,
            input [1:0][1:0] bin,
            output [1:0][1:0] cout);

sub1 U1 (.a(ain),
         .\b[1][1](bin[1][1]), .\b[1][0]),
         .\b[0][1](bin[0][1]), .\b[0][0]),
         .c(cout));
endmodule
\end{lstlisting}

% ============================================================
\section{允许引脚名同义词}
\subsection*{Allowing Pin Name Synonyms}

在 link process 中，可以根据指定的 synonym 启用 pin mapping 功能。该功能可解决下列各类不匹配：
\begin{itemize}
  \item 顶层例化与 RTL 块
  \item 逻辑库与 RTL 例化
  \item 逻辑库与物理库（FRAM 视图）
\end{itemize}

要启用引脚名同义词匹配：
\begin{enumerate}
  \item 确保 \varname{link\_allow\_pin\_name\_synonym} 变量设为 \texttt{true}（默认值）。
  \item 用 \cmd{set\_pin\_name\_synonym} 命令为 RTL、逻辑库与物理库引脚名指定同义词。
\end{enumerate}

要移除引脚名同义词，使用 \cmd{remove\_pin\_name\_synonym} 命令。
应在 link process 之后移除所有 pin name synonym，然后在加载 SDC 文件之前重新应用这些 synonym。
要报告引脚名同义词，使用 \cmd{report\_pin\_name\_synonym} 命令。

\subsection{示例 1}
\subsubsection*{Example 1}

下列示例在顶层例化 \texttt{top} 与 RTL 块 \texttt{sub} 之间启用引脚映射。
\cmd{set\_pin\_name\_synonym} 命令允许用 \texttt{s} 作为 \texttt{sel} 引脚的同义词，用 \texttt{do} 作为 \texttt{DOUT} 引脚的同义词。
\begin{lstlisting}
de_shell> set_app_var link_allow_pin_name_synonym true
de_shell> set_pin_name_synonym s sel
de_shell> set_pin_name_synonym do DOUT
\end{lstlisting}

\figplaceholder{Figure 12: Pin Mapping Between Top-Level Instantiations and RTL Blocks}{顶层例化与 RTL 块之间的引脚映射}{fig:pin-map-rtl}

\subsection{示例 2}
\subsubsection*{Example 2}

要在逻辑库与物理库之间启用引脚映射（如下例所示），\cmd{set\_pin\_name\_synonym} 命令：
\begin{enumerate}
  \item 创建 FRAM 视图，并保存在 Milkyway 设计库中。
  \item 根据指定的引脚名同义词，自动更正 FRAM 视图的引脚名以匹配逻辑库引脚名。
\end{enumerate}

该示例对 \texttt{do} 引脚使用同义词 \texttt{Dout}，对 \texttt{di} 引脚使用同义词 \texttt{Din}。
\begin{lstlisting}
de_shell> set link_allow_pin_name_synonym true
de_shell> set_pin_name_synonym Dout do
de_shell> set_pin_name_synonym Din di
\end{lstlisting}

\figplaceholder{Figure 13: Example of Pin Mapping Between Logic Library and Physical Library}{逻辑库与物理库之间引脚映射示例}{fig:pin-map-phys}

% ============================================================
\section{允许未连接的接口引脚}
\subsection*{Allowing Unconnected Interface Pins}

为防止 unconnected interface pin：
\begin{itemize}
  \item 连接到 constant high 或 constant low
  \item 跨 unconnected interface pin 传播 constant value
\end{itemize}
可将 \varname{link\_preserve\_dangling\_pins} 变量从默认值 \texttt{false} 改为 \texttt{true}。
将该变量设为 \texttt{true} 后，DC Explorer 可为每个 dangling interface pin 添加 anchor cell。
如图~\ref{fig:unconn-pins} 所示，DC Explorer 为未连接的 \texttt{wr} 信号添加了 \texttt{SNPS\_ANCHOR\_O} 锚点单元。

\figplaceholder{Figure 14: Example of Unconnected Interface Pins}{未连接接口引脚示例}{fig:unconn-pins}

% ============================================================
\section{在网表中包含锚点单元}
\subsection*{Including Anchor Cells in Netlists}

默认情况下，DC Explorer 会为 logic library 或 RTL 中的每个缺失 pin 创建 anchor cell，以防止该 pin 被移除并保留连接。
保存设计时，anchor cell 会自动从 netlist 中移除。
要在 Verilog 格式的 netlist 中包含 anchor cell，请对 \cmd{write\_file} 命令使用 \opt{-include\_anchor\_cells} 选项。

例如，下列命令以 Verilog 格式保存带锚点单元的层次设计：
\begin{lstlisting}
de_shell> write_file -hierarchy -format verilog \
   -output mapped.vg -include_anchor_cells
\end{lstlisting}

图~\ref{fig:anchor-netlist} 说明网表如何包含锚点单元。

\figplaceholder{Figure 15: Anchor Cells Support in Netlist}{网表中的锚点单元支持}{fig:anchor-netlist}

% ============================================================
\section{创建报告与缺失约束}
\subsection*{Creating Reports and Missing Constraints}

设计中可能包含不完整或不匹配的数据，例如 pin name 大小写不匹配、logic library 中缺失 cell、缺失 RTL module 以及 missing constraint。
不完整或缺失的 timing constraint 可能导致 infeasible path。
要报告不完整数据和 infeasible path，并创建 timing exception 和 missing constraint，请参阅：
\begin{itemize}
  \item 报告不完整或不匹配数据
  \item 报告缺失约束
  \item 报告不可行路径
  \item 使用分类时序报告创建约束
  \item 为不可行路径创建例外
\end{itemize}

\begin{seeAlsoBox}
容错类别。
\end{seeAlsoBox}

\subsection{报告不完整或不匹配数据}
\subsubsection*{Reporting Incomplete or Mismatched Data}

要报告不完整或不匹配数据，可任选下列方法之一：
\begin{itemize}
  \item 使用 \cmd{link} 命令显示所发现并已解析的不匹配警告列表。例如：
\begin{lstlisting}
Warning: Automatically reconstructed 'sel' bus with width of 2 from
{sel[0] ... sel[1]} bit-blasted bus pins in 't1' reference.
(LINK-912)

Warning: Automatically reconstructed 'd2' bus with width of 4 from
{ d2[0] d2[1] ... d2[3] } bit-blasted bus pins in 'U1' cell.
(LINK-912)

Warning: 'wr' pin on 'inst1' cell in 'top' design is missing from the
'SliceInit' reference. (LINK-902)

Warning: Unable to connect 'wr' net to the 'wr' pin of 'inst1' cell
in 'top' design. (LINK-903) . . . .

Warning: Design mismatches were detected by the linker and resolved
to link the current design. You can use report_design_mismatch
and check_design commands explicitly to see warning messages about
the design mismatches. (DESH-017)
\end{lstlisting}

  \item 使用 \cmd{check\_design} 命令显示工具遇到并解析的全部设计不匹配的摘要。例如：
\begin{lstlisting}
Design Mismatches                               Logical  Physical
------------------------------------------------------------------
Missing pin on logical library part (LINK-902)      1
Can't connect net to missing pin (LINK-903)         2
Missing physical pin (LINK-905)                              1
Presence of linker anchor cells (LINK-911)          2
Reconstructed bus (LINK-912)                        1
------------------------------------------------------------------
Total netlist elements                              6        1
\end{lstlisting}

  \item 使用 \cmd{report\_design\_mismatch} 命令显示设计不匹配的详细信息。例如：
\begin{lstlisting}
****************************************
Report:      design mismatches
Design:      mydesign
...
****************************************
Number of missing pins on logical library cells : 1

Design Object  Type    Mismatch
----------------------------------------------------------------
top    inst1   cell    'wr' pin on 'inst1' cell in 'top' design
                       is missing from the 'SliceInit' reference.

Number of unconnectable nets due to missing pins : 1

Design Object  Type    Mismatch
----------------------------------------------------------------
top    wr      net     Unable to connect 'wr' net to the
                       '*cell*6/wr' pin of 'inst1' cell in 'top'
                       design.

Number of linker anchor cells : 1

Design Object  Type    Mismatch
---------------------------------------------------------------
top    DCLINKER cell   Linker anchor cell inserted for missing
                       pin on 'inst1' cell.

Number of reconstructed bus : 1

Design Object  Type    Mismatch
---------------------------------------------------------------
top    U1      cell    Automatically reconstructed 'sel' bus
                       with width of 2 from { sel[0] ... sel[1]}
                       bit-blasted bus pins in 'my_design' reference.
\end{lstlisting}
\end{itemize}

\begin{seeAlsoBox}
\begin{itemize}
  \item 报告缺失约束
  \item 报告不可行路径
  \item 使用分类时序报告创建约束
\end{itemize}
\end{seeAlsoBox}

\subsection{报告缺失约束}
\subsubsection*{Reporting Missing Constraints}

要报告当前设计中缺失的约束（例如缺失时钟、I/O latency、驱动单元与输出负载），使用 \cmd{report\_missing\_constraints} 命令。例如：
\begin{lstlisting}
de_shell> report_missing_constraints
Information: Checking out the license 'DesignWare'. (SEC-104)
**************************************************
Report : Missing Constraints
Design : des_top
...
**************************************************
Missing Clock
----------------------------------------------------------------------
        clk_hi_1
----------------------------------------------------------------------
Total Missing Clock Definition found: 2.

Missing Input Delay
----------------------------------------
        combi_hi[3]
        comb1_hi[2]
----------------------------------------
The total number of missing input delay found is 2.

Missing Output Delay
----------------------------------------
        out_filter_i1
----------------------------------------
The total number of missing output delay found is 1.
\end{lstlisting}

默认情况下，DC Explorer 对缺失的 I/O delay 使用零值；对于缺失的 driving cell 和 output load，则从 logic library 中选择 inverter。
可以为 missing constraint 指定具体数值，而不采用工具推定的 constraint value。
要在 current design 中创建 missing constraint，请使用带相应选项的 \cmd{create\_missing\_constraints} 命令。

例如：
\begin{lstlisting}
de_shell> create_missing_constraints -period 0.75 \
   -input_delay 0.02 -output_delay 0.1 -output new1.sdc
\end{lstlisting}

\begin{seeAlsoBox}
\begin{itemize}
  \item 报告不可行路径
  \item 使用分类时序报告创建约束
  \item 为不可行路径创建例外
\end{itemize}
\end{seeAlsoBox}

\subsection{报告不可行路径}
\subsubsection*{Reporting Infeasible Paths}

Infeasible timing path 可能由下列原因造成：
\begin{itemize}
  \item 缺少 boundary condition
  \item input delay 的 timing goal 不切实际
  \item output delay 的 timing goal 不切实际
\end{itemize}

优化期间，\cmd{compile\_exploration} 命令会自动检测设计中的 infeasible path。
为控制优化工作聚焦于其他 path，该命令会暂时将这些 infeasible path 设置为 false path。
优化结束时，这些 false path 设置会被移除。

运行 \cmd{compile\_exploration} 之后，可以报告暂时被设为 false path 的 infeasible path。
为此，对 \cmd{create\_qor\_snapshot} 与 \cmd{query\_qor\_snapshot} 命令指定 \opt{-infeasible\_paths} 选项，然后生成 HTML 格式的分类时序报告，如下列命令序列所示。
要设置要报告的不可行路径最大数量，对 \cmd{create\_qor\_snapshot} 指定 \opt{-max\_paths} 选项。
\begin{lstlisting}
de_shell> compile_exploration
de_shell> create_qor_snapshot -max_paths 30 -infeasible_paths -name snap1
de_shell> query_qor_snapshot -name snap1 -infeasible_paths \
-columns {path_group infeasible_paths startpoint endpoint wns zero_path}\
-output_file snap1.htm
de_shell> sh firefox snap1.htm
\end{lstlisting}

图~\ref{fig:infeas-html} 给出由该命令序列生成的 HTML 格式分类时序报告。
优化期间检测到的每条 infeasible path 都会在 Infeasible Path 列中标记为 YES。

\figplaceholder{Figure 16: Categorized Timing Report in HTML Format}{HTML 格式的分类时序报告}{fig:infeas-html}

也可通过对 \cmd{report\_timing} 使用 \opt{-attributes} 选项生成时序报告来查看不可行路径。
报告显示起点与终点标有 \texttt{inf} 属性的每条不可行路径。

\begin{seeAlsoBox}
\begin{itemize}
  \item 为不可行路径创建例外
  \item 使用分类时序报告创建约束
  \item 分析结果质量
\end{itemize}
\end{seeAlsoBox}

\subsection{使用分类时序报告创建约束}
\subsubsection*{Creating Constraints Using Categorized Timing Reports}

可用 HTML 格式的分类时序报告创建下列约束：
\begin{itemize}
  \item False paths
  \item Multicycle paths
  \item Maximum delays
  \item Input delays
  \item Output delays
\end{itemize}

\begin{noteBox}
必须先用类似下列命令序列生成分类时序报告：
\begin{lstlisting}
de_shell> compile_exploration -scan
de_shell> create_qor_snapshot -nworst 3 -name my_snapshot
de_shell> query_qor_snapshot -name my_snapshot -output_file qor_report
de_shell> sh firefox qor_report.html
\end{lstlisting}
\end{noteBox}

\figplaceholder{Figure 17: Categorized Timing Report in HTML Format}{HTML 格式的分类时序报告}{fig:qor-html-create}

从图~\ref{fig:qor-html-create} 所示的分类时序 HTML 报告中，可按下列步骤创建缺失约束与例外：
\begin{enumerate}
  \item 单击 Append 按钮。
  \item 单击相应参数按钮，例如 Multicycle Paths。
  \item 为该参数输入数值。
  \item 单击 Create Exceptions 按钮。\\
        随即出现新的浏览器窗口，其中包含基于您所指定约束的一组命令。
  \item 用这些命令更新您的 SDC 文件。
\end{enumerate}

\begin{seeAlsoBox}
\begin{itemize}
  \item 度量结果质量
  \item 报告不可行路径
\end{itemize}
\end{seeAlsoBox}

\subsection{为不可行路径创建例外}
\subsubsection*{Creating Exceptions for Infeasible Paths}

可用 HTML 格式的分类时序报告为不可行路径创建例外。
必须先用类似下列命令序列生成包含设计中不可行路径的分类时序报告：
\begin{lstlisting}
de_shell> compile_exploration
de_shell> create_qor_snapshot -max_paths 30 -infeasible_paths -name snap1
de_shell> query_qor_snapshot -infeasible_paths -name snap1 \
-columns {path_group infeasible_paths startpoint endpoint wns zero_path}\
-output_file snap1.html
de_shell> sh firefox snap1.html
\end{lstlisting}

图~\ref{fig:infeas-exc} 给出由该命令序列生成的 HTML 格式分类时序报告。

\figplaceholder{Figure 18: Categorized Timing Report in HTML Format}{HTML 格式的分类时序报告}{fig:infeas-exc}

从分类时序 HTML 报告中，可按下列步骤创建缺失约束与例外：
\begin{enumerate}
  \item 单击 Append 按钮。
  \item 单击 False Paths 按钮。
  \item 在要为其创建例外的每条不可行路径的 Exceptions 列下勾选复选框。
  \item 单击 Create Exceptions 按钮。\\
        随即出现新的浏览器窗口，其中包含基于您所指定例外的一组命令。
        例如，当选择前五条不可行路径时，弹出窗口中显示下列命令：
\begin{lstlisting}
set_false_path -from u_clk2/X_out_reg/CLK -to out3;
set_false_path -from in1 -to u_in1/A_out_reg/D;
set_false_path -from in1 -to u_in2/A_out_reg/D;
set_false_path -from in3 -to u_clk2/A_out_reg/D;
set_false_path -from u_out3/X_out_reg/CLK -to out1;
\end{lstlisting}
  \item 用这些命令更新您的 SDC 文件。
\end{enumerate}

\begin{seeAlsoBox}
\begin{itemize}
  \item 报告不可行路径
  \item 报告结果质量
\end{itemize}
\end{seeAlsoBox}
